%% =====================================================================
%%  Chapitre 1 — Hello, Python!
%%  Tout le code est vérifié par outils/verifier_code.py (Shell IDLE simulé).
%% =====================================================================
\chapitre{1}{Hello, Python!}{Install Python and run your first programs}{illusHello}

\histoire{Every app on your phone, every website and every video game is a \textbf{program}: a
list of instructions that a computer follows, one after the other, exactly as written. Learning to
program means learning to write those instructions yourself. In this chapter, you install Python,
have your first conversation with it, then write, save and run your first real program. Along the
way, you meet Python's error messages\,---\,and discover that they are on your side.}

\begin{objectifs}
\item explain what a program, a programming language and an interpreter are;
\item install Python and open IDLE, the editor that comes with it;
\item type instructions in the Python shell and read its answers;
\item display text and numbers with \py{print()};
\item write, save and run a program file;
\item read an error message and fix the mistake it points to;
\item explain your code with comments.
\end{objectifs}

\auprogramme
  {what a program is; the interpreter; the shell and program files; errors}
  {\py{print()}, text between quotes, numbers, comments with \py{\#}}
  {IDLE: the Python shell and the code editor; keyboard shortcuts}
  {read error messages from the bottom up; save and run often}
  {a business card that Python prints for you}

\seances[4]{livre/otNavy/Get ready/What programming is and how to install Python (1.1--1.2),
  terminal/otViolet/Talk to Python/The shell and the \py{print()} function (1.3--1.4),
  fichier/otTeal/Your first file/Programs and error messages (1.5--1.8),
  crayon/otGreenDark/Practise/Mini-project and exercises and review (1.9--1.11)}

%% ---------------------------------------------------------------------
\rubrique{learn}{What is a program?}

A computer is an extraordinary machine: it performs billions of operations per second, never gets
tired and never forgets anything. But it has one weakness: it does \emph{nothing} on its own. Every
single thing it does, it does because a program told it to.\otPara

A program is a bit like a cooking recipe. A recipe is a list of steps\,---\,break two eggs, add
100~g of flour, stir for one minute\,---\,that you follow in order, from the first to the last. A
program is a list of \textbf{instructions} that the computer carries out in order, from the first to
the last. There is one big difference: a cook has common sense, a computer has none. If a recipe
says \enquote{add some salt}, a cook guesses how much. A computer cannot guess: every instruction must
be perfectly precise.

\begin{definition}{program}
A list of instructions that a computer carries out, one after the other, to perform a task. Each
instruction is a simple, precise order: display a message, add two numbers, save a file\ldots
\end{definition}

\soustitre{Talking to a computer}
Deep inside, a computer only understands long sequences of 0s and 1s, which are almost impossible
for humans to read. Human languages, on the other hand, are full of ambiguities that a computer
cannot handle. A \textbf{programming language} sits in between: precise enough for the computer,
readable enough for you. There are hundreds of them\,---\,JavaScript, Java, C, Swift\ldots\,---\,and
Python is one of the most popular in the world.\otPara

Here is what a short Python program looks like. You don't need to understand it yet: just notice
how close it is to plain English.

\begin{code}{file=first_look.py}
print("Hello!")
print("Python can count:", 1 + 2 + 3)
print("Goodbye!")
\end{code}
\begin{sortie}
Hello!
Python can count: 6
Goodbye!
\end{sortie}

Each line is one instruction. The word \py{print} means \enquote{display on the screen}. The
computer carried out the three instructions in order, from top to bottom, and even did the
addition for us. That's a program!

\begin{definition}{programming language}
A language with strict rules, designed to write instructions that a computer can carry out. The
rules that say how instructions must be written are called the \textbf{syntax} of the language.
\end{definition}

\soustitre{The interpreter, Python's translator}
When you run a Python program, a piece of software called the Python \textbf{interpreter} reads your
instructions one by one, checks that they follow the rules of the language, and carries them out.
When people say \enquote{install Python}, they mean installing this interpreter on their computer.

\begin{center}
\begin{tikzpicture}[x=1mm,y=1mm]
  \foreach \x/\ico/\coul/\titre/\txt in {0/clavier/otNavy/You/write instructions in Python,
      58/engrenage/otViolet/The interpreter/reads them and carries them out,
      116/ecran/otGreenDark/The computer/shows the result}{%
    \fill[\coul!8,rounded corners=2mm] (\x,0) rectangle ++(50,26);
    \fill[\coul,rounded corners=1.6mm] (\x+4,15) rectangle ++(8,8);
    \otDessinIcone{\ico}{white}{\coul}{5.2mm}{(\x+8,19)}
    \node[anchor=west,font=\otSemi\fontsize{9.6}{11}\selectfont,text=\coul] at (\x+14,19) {\titre};
    \node[anchor=north west,text width=42mm,font=\fontsize{9.4}{11.6}\selectfont,text=otInk] at (\x+3.5,12) {\txt};}
  \foreach \x in {50.5,108.5}{\draw[-{Stealth[length=2.6mm]},otMuted,line width=1pt] (\x,13) -- ++(7,0);}
\end{tikzpicture}
\end{center}

\begin{definition}{interpreter}
The program that reads Python code and carries it out, instruction by instruction. Installing
\enquote{Python} on a computer means installing the Python interpreter.
\end{definition}

\begin{culture}[Where does the name come from?]
Python was created by the Dutch programmer Guido van Rossum, who released the first version in
1991. The name has nothing to do with snakes: it comes from \emph{Monty Python's Flying Circus}, a
British comedy series he enjoyed. Today Python is free, open source and used almost everywhere:
websites, scientific research, data analysis, artificial intelligence, video games, and the
automation of countless everyday tasks.
\end{culture}

\soustitre{Why start with Python?}
\begin{itemize}
\item \textbf{It reads almost like English.} You can focus on ideas instead of punctuation.
\item \textbf{It is versatile.} The same language lets you automate a spreadsheet, build a website
  or train an artificial intelligence.
\item \textbf{It is free} and works the same way on Windows, macOS and Linux.
\item \textbf{It has a huge community.} Whatever your question, someone has probably answered it
  already, and thousands of free libraries extend what Python can do.
\end{itemize}

%% ---------------------------------------------------------------------
\rubrique{learn}{Installing Python}

Python is free. Installing it takes about ten minutes, and you only do it once. Follow the steps for
your system.

\begin{encadre}{ecran}{Windows}[]
\begin{etapes}
\item Open your web browser and go to \textbf{python.org}. Click \menu{Downloads}: the site detects
  that you use Windows and offers the latest version of Python~3.
\item Download it, then open the downloaded file.
\item Follow the instructions and keep the default options. If the installer shows a box labelled
  \menu{Add python.exe to PATH}, tick it: it will be useful later.
\item When the installation is over, open the \menu{Start} menu and type \menu{IDLE}. Click the
  application \menu{IDLE (Python 3.x)}.
\end{etapes}
\end{encadre}

\begin{encadre}{ecran}{macOS}[]
\begin{etapes}
\item Go to \textbf{python.org}, click \menu{Downloads} and download the latest version of Python~3
  for macOS.
\item Open the downloaded \textbf{.pkg} file and follow the instructions.
\item Open your \menu{Applications} folder, then the \menu{Python 3.x} folder, and double-click
  \menu{IDLE}.
\end{etapes}
\end{encadre}

\begin{encadre}{ecran}{Linux}[]
Python~3 is almost always installed already, but IDLE sometimes isn't. On Ubuntu or Debian, open a
terminal and type \py{sudo apt install idle3}, then start IDLE from your applications menu, or by
typing \py{idle3} in the terminal. On other distributions, look for the \enquote{idle} package in
your package manager.
\end{encadre}

\begin{astuce}[Which version?]
Any version of Python~3 from 3.10 onwards works with this book. Every program in it was run with
Python~3.13; if your version is more recent, a few messages may be worded slightly differently.
Never install Python~2: it has been obsolete since 2020.
\end{astuce}

\soustitre{Meet IDLE}
IDLE is the small programming environment installed with Python. Its name stands for
\emph{Integrated Development and Learning Environment}\,---\,and it is also a wink at Eric Idle, one
of the Monty Python team. When you open it, you see a window like this one:

\begin{idleshell}{IDLE Shell 3.13.5}
Python 3.13.5 (...) on win32
Enter "help" below or click "Help" above for more information.
>>>
\end{idleshell}

The first line shows the version of Python (3.13.5 here; yours may be more recent) and your system
(\py{win32} for Windows, \py{darwin} for macOS, \py{linux} for Linux). The three chevrons
\py{>>>} are the \textbf{prompt}: Python is waiting for your instructions.

\begin{definition}{IDLE}
The editor installed with Python. It has two kinds of windows: the \textbf{Shell}, to talk to Python
one instruction at a time, and the \textbf{editor}, to write and save complete programs.
\end{definition}

%% ---------------------------------------------------------------------
\rubrique{learn}{Your first conversation: the shell}

Click in the IDLE Shell window, after the prompt \py{>>>}. Type \py{2 + 3}, then press
\touche{Enter}. Python answers at once:

\begin{console}
>>> 2 + 3
5
\end{console}

You have just run your first instruction! Python read \py{2 + 3}, worked out the result and
displayed it on the next line. Then it showed a new prompt, ready for the next instruction. Try a
few more calculations; the star \py{*} means \enquote{multiplied by}:

\begin{console}
>>> 10 - 4
6
>>> 7 * 8
56
>>> 1 + 2 + 3 + 4 + 5
15
\end{console}

\begin{regle}[How to read this book]
In this book, a dark box labelled \textbf{Python shell} shows a conversation with Python. The lines
that start with \py{>>>} are what \emph{you} type; the other lines are Python's answers. Type only
what comes after \py{>>>}: the chevrons are already on your screen.
\end{regle}

\soustitre{Your first words}
Now let's display some text. Type the following instruction exactly as it is written, with the
parentheses and the quotation marks:

\begin{console}
>>> print("Hello, world!")
Hello, world!
\end{console}

Python displayed the text that was between the quotation marks. \enquote{Hello, world!} is a
tradition: it is the first program of millions of programmers around the world.

\begin{culture}[Hello, world!]
The tradition of starting with a program that displays \enquote{Hello, world!} was made famous in
1978 by the book \emph{The C Programming Language}, by Brian Kernighan and Dennis Ritchie. Almost
every programming course since has followed it\,---\,this one included.
\end{culture}

The shell is perfect for trying things out: you type, Python answers immediately. It has one
limit: it forgets everything when you close it. To keep a program, you will write it in a file
(section~1.5).

\begin{atoi}[Calculate in the shell]
In the shell, work out \py{15 + 27}, \py{100 - 58} and \py{6 * 7}. Then make Python display
\py{Hello, Python!} with \py{print()}. What do you notice about the three results?
\end{atoi}

%% ---------------------------------------------------------------------
\rubrique{learn}{Displaying with \texorpdfstring{\py{print()}}{print()}}

\py{print} is a \textbf{function}: a ready-made tool that Python offers you. To use it, you write
its name followed by parentheses, and you put inside the parentheses what you want to display.
Using a function is called \textbf{calling} it.

\begin{definition}{function}
A named tool that performs a task. You call a function by writing its name followed by
parentheses: \py{print("Hi")}. The values inside the parentheses are its \textbf{arguments}.
\end{definition}

\soustitre{Text goes between quotes}
In Python, a piece of text is called a \textbf{string}. You write it between quotation marks,
either double quotes \py{"} or single quotes \py{'}. Both work the same way, as long as you use the
same kind at both ends:

\begin{console}
>>> print("I am learning Python.")
I am learning Python.
>>> print('Me too!')
Me too!
\end{console}

Notice that the quotes are not displayed: they only tell Python where the text begins and where it
ends. Inside a string, you can write anything: letters, spaces, digits, punctuation.

\soustitre{Numbers go without quotes}
Numbers are written without quotes, just as in mathematics. Decimal numbers use a \textbf{dot},
never a comma:

\begin{console}
>>> print(42)
42
>>> print(3.5)
3.5
>>> print(7 * 6)
42
\end{console}

Now compare these two instructions carefully:

\begin{console}
>>> print(2 + 3)
5
>>> print("2 + 3")
2 + 3
\end{console}

Without quotes, \py{2 + 3} is a calculation: Python works it out and displays \py{5}. With quotes,
\py{"2 + 3"} is just text: Python displays it exactly as it is, character by character.

\begin{regle}[Quotes or no quotes?]
Between quotes, Python takes your text \emph{literally} and does not try to understand it. Without
quotes, Python reads and \emph{evaluates} what you wrote: a calculation is computed, a name is looked
up.
\end{regle}

\soustitre{Several values at once}
You can give \py{print()} several values, separated by commas. Python displays them on the same
line, with a space between two values:

\begin{console}
>>> print("Total:", 2 + 3)
Total: 5
>>> print("I am", 16, "years old")
I am 16 years old
\end{console}

Python added the spaces itself: there is no need to put one at the end of \py{"I am"}. Finally,
\py{print()} with nothing between the parentheses displays an empty line, which is useful to space
out what your program displays.

\begin{atoi}[Predict, then check]
Before typing them, predict what each instruction displays. Then check in the shell.\par\vspace{1mm}
\py{print("10 + 5")}\qquad\py{print(10 + 5)}\qquad\py{print("10", "+", 5)}\qquad\py{print(10, 5)}
\end{atoi}

%% ---------------------------------------------------------------------
\rubrique{learn}{Your first program file}

The shell forgets everything as soon as you close it. A real program is saved in a \textbf{file},
so that you can run it again whenever you want, improve it and share it. A Python program file is
called a \textbf{script}; its name ends with \textbf{.py}.

\begin{etapes}
\item In IDLE, choose \menu{File\menusep New File} (or press \touche{Ctrl}\,+\,\touche{N}; on a Mac,
  \touche{Cmd}\,+\,\touche{N}). An empty editor window opens.
\item Type your program in this window. Nothing happens while you type: the editor just records your
  instructions.
\item Save it with \menu{File\menusep Save} (\touche{Ctrl}\,+\,\touche{S}). Create a folder called
  \textbf{python} in your \textbf{Documents} folder, and save the file as \textbf{hello.py}.
\item Run it with \menu{Run\menusep Run Module}, or simply press \touche{F5}. The result appears in
  the Shell window.
\end{etapes}

Here is your first program, as it looks in the IDLE editor. Notice that IDLE colours the code to
make it easier to read: functions in purple, text in green.

\begin{idleediteur}{hello.py - C:\textbackslash Users\textbackslash ana\textbackslash Documents\textbackslash python\textbackslash hello.py (3.13.5)}
print("Hello, world!")
print("My name is Ana.")
print("I am learning Python.")
\end{idleediteur}

Press \touche{F5}. In the Shell window, IDLE first writes a line \py{= RESTART:} followed by the name
of your file (the Shell starts afresh to run it), then the three lines of your program appear:

\rattacher
\titresortie{IDLE Shell}
\begin{sortie}
Hello, world!
My name is Ana.
I am learning Python.
\end{sortie}

\begin{regle}[How to read this book (continued)]
From now on, programs are shown in a light box like the one below, with the name of the file at the
top and line numbers in the margin (do not type the numbers!). What the program displays when you
run it is shown just below, in a dark box labelled \textbf{Output}.
\end{regle}

\begin{code}{file=morning.py}
print("Good morning!")
print("Today is a great day to learn Python.")
print()
print("Lesson 1: the print() function")
\end{code}
\begin{sortie}
Good morning!
Today is a great day to learn Python.

Lesson 1: the print() function
\end{sortie}

\soustitre{From top to bottom}
Python runs the instructions of a script in order, starting with line~1 and ending with the last
line. This is called \textbf{sequential execution}. If you swap two lines, you change the order of
what is displayed: the computer never reorders anything by itself.

\begin{astuce}[Save, run, repeat]
Get into the habit of pressing \touche{F5} after every small change: it saves your file and runs
it. Programmers rarely write a whole program in one go; they add a few lines, run, check, and start
again. The sooner you see a mistake, the easier it is to fix.
\end{astuce}

\begin{atoi}[About me]
Create a new file \textbf{about\_me.py} that displays three lines: your first name, the city where
you live and your favourite food. Save it and run it with \touche{F5}.
\end{atoi}

%% ---------------------------------------------------------------------
\rubrique{hood}{What happens when you press F5?}

Pressing \touche{F5} sets off a short chain of events. Understanding it will help you make sense of
the errors you meet.

\begin{center}
\begin{tikzpicture}[x=1mm,y=1mm]
  \foreach \k/\ico/\titre/\txt in {0/fichier/Save/IDLE saves your file,
      1/boucle/Restart/a fresh Python starts,
      2/oeil/Check/the whole file is read and its syntax checked,
      3/engrenage/Translate/into simpler instructions (bytecode),
      4/lecture/Run/line by line from top to bottom}{%
    \pgfmathsetmacro{\x}{\k*34.4}
    \fill[otViolet!9,rounded corners=2mm] (\x,0) rectangle ++(31,30);
    \fill[otViolet,rounded corners=1.6mm] (\x+11.5,19) rectangle ++(8,8);
    \otDessinIcone{\ico}{white}{otViolet}{5.2mm}{(\x+15.5,23)}
    \node[font=\otSemi\fontsize{9}{10}\selectfont,text=otViolet] at (\x+15.5,14.5) {\titre};
    \node[anchor=north,text width=28mm,align=center,font=\fontsize{8.4}{10}\selectfont,text=otInk] at (\x+15.5,11) {\txt};
    \ifnum\k<4\draw[-{Stealth[length=2mm]},otMuted,line width=0.9pt] (\x+31.6,15) -- ++(2.4,0);\fi}
\end{tikzpicture}
\end{center}

The important point is step~3: \textbf{Python reads and checks your whole file before running the
first line.} If it finds an instruction that breaks the rules of the language\,---\,a
\textbf{syntax error}\,---\,it refuses to run anything at all. Look at this program, whose last line
has lost its closing parenthesis:

\begin{code}{file=check.py, bug}
print("Line 1")
print("Line 2")
print("Line 3"
\end{code}
\begin{erreur}
SyntaxError: '(' was never closed
\end{erreur}

Not even \enquote{Line 1} is displayed! IDLE shows this message in a small window titled
\textbf{SyntaxError}, and highlights in red the place where it got lost in your code. Now compare
with this program, which is written correctly but uses a function that does not exist, \py{prin}:

\begin{code}{file=check.py, bug}
print("Line 1")
print("Line 2")
prin("Line 3")
print("Line 4")
\end{code}
\begin{erreur}
Line 1
Line 2
Traceback (most recent call last):
  File "check.py", line 3, in <module>
    prin("Line 3")
NameError: name 'prin' is not defined. Did you mean: 'print'?
\end{erreur}

This time, the syntax is correct, so the program starts. Lines~1 and~2 run normally. Python only
discovers the problem when it reaches line~3: it stops there, and line~4 never runs.

\begin{regle}[Two kinds of errors]
A \textbf{syntax error} is detected before the program starts: nothing runs. An error that happens
\textbf{while the program runs} (Python calls it an \textbf{exception}) stops the program at the
line where it occurs; everything before that line has already run.
\end{regle}

\begin{culture}[Bytecode]
Step~4 is invisible: Python translates your code into \emph{bytecode}, a list of very simple
instructions that its interpreter can run quickly. You will never need to read bytecode, but it
explains why Python is called an \enquote{interpreted} language.
\end{culture}

%% ---------------------------------------------------------------------
\rubrique{debug}{Reading error messages}

Every programmer, beginner or expert, makes mistakes all day long. Professionals are not people who
never make mistakes: they are people who know how to \emph{read} the messages Python gives them.
These messages are not reproaches; they are clues.

\soustitre{The anatomy of an error message}
Here is the message from the previous section again. Each part tells you something:

\sortielibre
\begin{erreur}
Traceback (most recent call last):  <<1>>
  File "check.py", line 3, in <module>  <<2>>
    prin("Line 3")  <<3>>
NameError: name 'prin' is not defined. Did you mean: 'print'?  <<4>>
\end{erreur}

\begin{renvois}
\renvoi{1} \textbf{Traceback} announces an error message: Python is going to retrace the steps that
  led to the error.
\renvoi{2} \textbf{Where}: the file (\py{check.py}) and the line number (\py{line 3}). \py{<module>}
  means \enquote{the main program}. In IDLE, the full path of the file is shown, for example
  \py{C:\textbackslash Users\textbackslash ana\textbackslash Documents\textbackslash python\textbackslash check.py};
  in this book we shorten it to the name of the file.
\renvoi{3} \textbf{The faulty line}, copied from your program.
\renvoi{4} \textbf{What went wrong}: the type of error (\py{NameError}), followed by an explanation
  in plain English, and sometimes a suggestion. This is the most important line.
\end{renvois}

\begin{regle}[Read from the bottom up]
Always start with the \textbf{last line}: it tells you \emph{what} is wrong. Then go up to find
\emph{where}: the line number takes you straight to the problem.
\end{regle}

\soustitre{The six most common beginner mistakes}
Try each of these wrong instructions in the shell: you will recognise these messages when they
appear in your own programs. In the shell, a syntax error is reported with a single line.

\cas{1}{A capital letter}
Python is \textbf{case-sensitive}: for it, \py{Print} and \py{print} are two different names, and
only \py{print} exists. The \py{<pyshell\#0>} part means \enquote{the instruction you typed in the
shell}.
\begin{console}
>>> Print("Hello")
Traceback (most recent call last):
  File "<pyshell#0>", line 1, in <module>
    Print("Hello")
NameError: name 'Print' is not defined. Did you mean: 'print'?
\end{console}

\cas{2}{Forgotten quotes}
Without quotes, \py{Hello} is not text: Python looks for something named \py{Hello}, finds nothing
and gives up. Fix: \py{print("Hello")}.
\begin{console}
>>> print(Hello)
Traceback (most recent call last):
  File "<pyshell#0>", line 1, in <module>
    print(Hello)
NameError: name 'Hello' is not defined
\end{console}

\cas{3}{A missing closing quote}
The text never ends: in Python, a string must open and close on the same line. \enquote{String
literal} is the official name of a piece of text written in your code.
\begin{console}
>>> print("Hello)
SyntaxError: unterminated string literal (detected at line 1)
\end{console}

\cas{4}{Unbalanced parentheses}
Every opening parenthesis needs exactly one closing parenthesis. (In the IDLE Shell, if you forget
the closing parenthesis and press \touche{Enter}, IDLE simply waits for the rest of the instruction:
type \py{)} and press \touche{Enter} again.)
\begin{console}
>>> print("Hello"
SyntaxError: '(' was never closed
>>> print("Hello"))
SyntaxError: unmatched ')'
\end{console}

\cas{5}{No parentheses}
Very old versions of Python allowed this; Python~3 does not, and even suggests the fix.
\begin{console}
>>> print "Hello"
SyntaxError: Missing parentheses in call to 'print'. Did you mean print(...)?
\end{console}

\cas{6}{A space at the beginning of the line}
In Python, spaces at the start of a line (the \textbf{indentation}) have a meaning that you will
discover in chapter~8. For now, start every line at the very left.
\begin{console}
>>>  print("Hello")
SyntaxError: unexpected indent
\end{console}

\begin{pratique}[A method that always works]
\begin{etapes}
\item Read the \textbf{last line} of the message: what type of error, what explanation?
\item Go to the \textbf{line number} given in the message.
\item Look closely at that line\,---\,and at the end of the line before, for a forgotten quote or
  parenthesis.
\item Change \textbf{one thing}, then run the program again.
\end{etapes}
\end{pratique}

%% ---------------------------------------------------------------------
\rubrique{pro}{Comments and clean code}

Programs are read by computers, but also by people: your teammates, and yourself in three months'
time. Professional programmers write their code so that a human can read it easily.

\soustitre{Comments}
Everything that follows a hash sign \py{\#} on a line is a \textbf{comment}: Python ignores it
completely. Comments are notes for human readers.

\begin{code}{file=comments.py}
# comments.py - my first commented program
# Author: Ana, September 2026

print("Comments are ignored by Python.")  # this one too
# print("This line will not run.")
print("Only two lines are displayed.")
\end{code}
\begin{sortie}
Comments are ignored by Python.
Only two lines are displayed.
\end{sortie}

Line~5 shows a classic trick: putting a \py{\#} in front of an instruction \enquote{switches it
off} without deleting it. This is called \textbf{commenting out} a line. In IDLE, you can comment
out the selected lines with \menu{Format\menusep Comment Out Region} (\touche{Alt}\,+\,\touche{3}).

\begin{definition}{comment}
A note in the code, written after a \py{\#} sign. Python ignores it: it is there to help humans
understand the program.
\end{definition}

\begin{pratique}[Comment the why, not the what]
A good comment explains what the code cannot say by itself: the purpose of the program, a choice
you made, a trap to avoid. The comment \py{\# display Hello} above \py{print("Hello")} adds nothing.
Start each program with a comment that says what it does.
\end{pratique}

\soustitre{Clean code from day one}
Python's official style guide, called \textbf{PEP~8}, describes how to lay out code so that every
Python programmer can read it easily. Here are its first rules:

\begin{itemize}
\item one instruction per line;
\item no space before the opening parenthesis: \py{print("Hi")}, not \py{print ("Hi")};
\item a space after each comma, not before: \py{print("Total:", 5)};
\item file names in lower case, without spaces or accents, words joined by an underscore
  \py{\_}: \py{about\_me.py}, \py{business\_card.py}.
\end{itemize}

\begin{attention}[Do not use Python's own names]
Never call your file \py{print.py}, \py{random.py} or \py{turtle.py}: these names belong to Python
and its tools, and a file with the same name in your folder would get in the way. Choose names that
describe \emph{your} program.
\end{attention}

%% ---------------------------------------------------------------------
\rubrique{build}{Mini-project: your business card}

Time to put everything together. You are going to write a program that displays a business card
with a frame drawn with characters, like this:

\cible{carte}
\begin{sortie}
+--------------------------------+
|                                |
|   ADA LOVELACE                 |
|   Programmer in training       |
|                                |
|   ada@example.com              |
|   Learning Python since today  |
|                                |
+--------------------------------+
\end{sortie}

\soustitre{Step 1: plan the lines}
Before typing, count: the card has nine lines, so the program needs nine \py{print()}
instructions. The top and bottom borders are made of \py{+} and \py{-} signs; the other lines start
and end with a vertical bar \py{|}.

\soustitre{Step 2: write the borders and the text}
Every line of the card must have \textbf{exactly the same length}, here 34 characters, otherwise the
right-hand border will be jagged. The Shell uses a \textbf{monospaced} font, in which every
character, even a space, takes the same width: if two lines have the same number of characters,
they are exactly as long as each other. Fill the end of each line with spaces up to the bar.

\begin{code}{file=business_card.py, cible=carte}
# business_card.py - displays my business card
print("+--------------------------------+")
print("|                                |")
print("|   ADA LOVELACE                 |")
print("|   Programmer in training       |")
print("|                                |")
print("|   ada@example.com              |")
print("|   Learning Python since today  |")
print("|                                |")
print("+--------------------------------+")
\end{code}

\soustitre{Step 3: run and adjust}
Press \touche{F5}. If a bar on the right sticks out, count the characters of that line: one space
too many or too few is enough. IDLE helps you: at the bottom right of the editor window, it shows
the line (\textbf{Ln}) and the column (\textbf{Col}) where your cursor is.

\begin{culture}[Ada Lovelace]
Ada Lovelace (1815--1852) was an English mathematician. In 1843 she published notes on the
\emph{Analytical Engine}, a mechanical computer designed by Charles Babbage, including a method to
make it calculate a sequence of numbers. It is often described as the first computer program ever
published.
\end{culture}

\begin{atoi}[Make it yours]
Change the card so that it shows your own name and details. Then make the card two characters wider
(add a \py{-} to the borders and a space to every other line). Run it to check that the frame is
still perfectly straight.
\end{atoi}

%% ---------------------------------------------------------------------
\rubrique{practice}{Exercises}

\exercice[type=predict,niveau=1]{Text or calculation?}{For each instruction, write exactly what
Python displays. Then check your answers in the shell.}
\begin{questions}[2]
\item \py{print("Python")}
\item \py{print("4 * 5")}
\item \py{print(4 * 5)}
\item \py{print("4", "*", 5)}
\item \py{print(10 - 3, "days")}
\item \py{print("Total:", 10 + 15)}
\end{questions}

\exercice[type=predict,niveau=1]{In order, please}{Write exactly what this program displays,
line by line, including any empty line.}
\begin{code}{file=order.py}
print("3")
print("2")
print("1")
print()
print("Lift-off!", 3 + 2 + 1, "seconds later")
\end{code}
\ecranvide{5}

\exercice[type=shell,niveau=1]{A pocket calculator}{Use the shell to answer these questions. Write
the instruction you typed and Python's answer.}
\begin{questions}
\item A cinema sold 125 tickets in the morning and 389 in the afternoon. How many in total?
\item You had 1000 euros and spent 258. How much is left?
\item A box contains 12 rows of 12 eggs. How many eggs are there?
\item Type \py{2 + 3 * 4}. Is the result \py{20}? What can you deduce about the order in which
  Python calculates?
\end{questions}

\exercice[type=fix,niveau=1]{Spot the mistake}{Each instruction contains one mistake. For each one,
predict the type of error (\py{NameError} or \py{SyntaxError}), then write the correct
instruction.}
\begin{questions}[2]
\item \py{Print("Good morning")}
\item \py{print("Good morning)}
\item \py{print(Good morning)}
\item \py{print "Good morning"}
\item \py{pritn("Good morning")}
\item \py{print("Good morning"))}
\end{questions}

\exercice[type=write,niveau=1]{Hello from me}{Write a program \py{hello\_me.py} that displays the
three lines below, with your own details instead of Tom's.}
\cible{hello-me}
\begin{sortie}
My name is Tom.
I live in Manchester.
My favourite food is pizza.
\end{sortie}

\exercice[type=complete,niveau=2]{Fill in the blanks}{Complete the program so that it displays
exactly the output shown below. Each blank hides one word, one number or one piece of text.}
\begin{code}{file=blanks.py, trous}
print(________________)
print("Lunch is at", ____)
_____("See you later!")
print(_____, "+", 2, "=", 7)
\end{code}
\cible{blanks}
\begin{sortie}
Good morning!
Lunch is at 12
See you later!
5 + 2 = 7
\end{sortie}

\exercice[type=read,niveau=2]{Read the message}{Ana ran the program below and got this message.}
\begin{code}{file=menu.py, bug}
print("Today's menu")
print("Starter: soup")
print("Main course: fish")
prnit("Dessert: apple pie")
print("Enjoy your meal!")
\end{code}
\begin{erreur}
Today's menu
Starter: soup
Main course: fish
Traceback (most recent call last):
  File "menu.py", line 4, in <module>
    prnit("Dessert: apple pie")
NameError: name 'prnit' is not defined. Did you mean: 'print'?
\end{erreur}
\begin{questions}
\item What type of error is it? On which line is the mistake?
\item Why were the first three lines displayed, but not the last one?
\item Correct the program. What does it display now?
\end{questions}

\exercice[type=fix,niveau=2]{Debug the poem}{This program should display the short poem shown
below, but it contains three mistakes. Find them by running the program several times: each run
reveals one mistake. Write the corrected program.}
\begin{code}{file=poem.py, bug}
print("Roses are red,")
print("Violets are blue,)
Print("Python is fun,")
print("And so are you!"
\end{code}
\cible{poem}
\begin{sortie}
Roses are red,
Violets are blue,
Python is fun,
And so are you!
\end{sortie}

\exercice[type=write,niveau=2]{A tree in the shell}{Write a program \py{tree.py} that draws this
tree, with one \py{print()} per line. Then add a few decorations of your own.}
\cible{tree}
\begin{sortie}
     *
    ***
   *o***
  ***o***
 *o*****o*
    | |
  =======
\end{sortie}

\exercice[type=predict,niveau=2]{Ghost lines}{Some lines of this program are comments. Write what
the program displays, then explain why.}
\begin{code}{file=ghosts.py}
# ghosts.py - what will appear?
print("one")
# print("two")
print("three")  # print("four")
#print("five")
print("six # seven")
\end{code}
\ecranvide{3}

\exercice[type=write,niveau=2]{The 7 times table}{Write a program that displays the first five lines
of the 7 times table, as below. Python must do the multiplications: write \py{7 * 3}, not \py{21}.}
\cible{table7}
\begin{sortie}
7 x 1 = 7
7 x 2 = 14
7 x 3 = 21
7 x 4 = 28
7 x 5 = 35
\end{sortie}

\exercice[type=challenge,niveau=3]{The café receipt}{At the Corner Café, a coffee costs 3 euros, a
muffin 4 euros and an orange juice 5 euros. Lucy orders 2 coffees, 1 muffin and 3 orange juices.
Write a program that displays her receipt exactly as below. Python must calculate every amount and
the total itself: the only numbers you may type are the quantities and the prices.}
\cible{receipt}
\begin{sortie}
====== CORNER CAFE ======
2 x Coffee       = 6
1 x Muffin       = 4
3 x Orange juice = 15
-------------------------
TOTAL = 25 euros
Thank you, see you soon!
\end{sortie}

%% ---------------------------------------------------------------------
\rubrique{review}{Chapter review}

\begin{essentiel}
\otEssTitre{Programs and Python}
A \textbf{program} is a list of instructions that the computer carries out in order, from top to
bottom. The \textbf{interpreter} reads Python code and runs it. \textbf{IDLE} has two kinds of
windows: the \textbf{Shell} (\py{>>>}), for trying one instruction at a time, and the
\textbf{editor}, for writing programs saved in \textbf{.py} files and run with \touche{F5}.

\otEssTitre{Displaying with print()}
\begin{code}{mini}
print("Hello, world!")        # text: between quotes
print(42, 3.5)                # numbers: no quotes, decimal dot
print("Total:", 2 + 3)        # several values: Total: 5
print()                       # an empty line
\end{code}

\otEssTitre{Comments}
Everything after \py{\#} is ignored by Python. A comment explains \emph{why}; start every program
with one that says what it does.

\otEssTitre{Errors}
A \textbf{syntax error} (\py{SyntaxError}) stops everything before the first line runs. An
\textbf{exception} such as \py{NameError} stops the program at the faulty line. Read the message
\textbf{from the bottom up}: the last line says what, the line number says where.
\end{essentiel}

\begin{quiz}
\qcm{What is a program?}{a computer}{a list of instructions}{a kind of file}
\qcm{What does the interpreter do?}{it translates English into Python}{it reads Python code and
  runs it}{it draws the windows of IDLE}
\qcm{In the Shell, what does \py{>>>} mean?}{Python is waiting for you}{an error occurred}{the
  program is finished}
\qcm{What does \py{print("3 + 4")} display?}{\py{7}}{\py{3 + 4}}{\py{"3 + 4"}}
\qcm{What does \py{print("Age:", 12)} display?}{\py{Age:12}}{\py{Age: 12}}{\py{"Age:", 12}}
\qcm{Which instruction is correct?}{\py{Print("Hi")}}{\py{print("Hi")}}{\py{print "Hi"}}
\qcm{Which extension does a Python program file have?}{\py{.txt}}{\py{.python}}{\py{.py}}
\qcm{What does Python do with \py{\# Hello}?}{it ignores it}{it displays \py{Hello}}{it reports an
  error}
\qcm{A program has a syntax error on line~5. What runs?}{lines 1 to 4}{nothing at all}{every
  line except line~5}
\qcm{Where do you look first in an error message?}{the first line}{the last line}{the line
  number}
\end{quiz}

\begin{jesais}
\item explain what a program, a programming language and an interpreter are;
\item open IDLE and run instructions in the Shell;
\item display text and numbers, alone or together, with \py{print()};
\item write a program in the editor, save it in a \textbf{.py} file and run it with \touche{F5};
\item read an error message from the bottom up and find the faulty line;
\item explain my code with comments, and switch a line off by commenting it out.
\end{jesais}
